\section{Diagnóstico}

\graphicspath{{../trabajos/07_figuras_para_icr/planos_qgis/}{../trabajos/07_figuras_para_icr/fotografias_icr_optimizadas/C-01_penca_de_maguey/}{../trabajos/07_figuras_para_icr/fotografias_icr_optimizadas/C-02_piedra_local_tepetate/}{../trabajos/07_figuras_para_icr/fotografias_icr_optimizadas/C-03_bajareque/}{../trabajos/07_figuras_para_icr/fotografias_icr_optimizadas/C-04_organos/}}

Este capítulo constituye el núcleo empírico de la ICR. Su objetivo es ordenar la evidencia de campo para mostrar qué sistemas constructivos tradicionales persisten en las comunidades estudiadas, cómo se distribuyen en el territorio y qué cambios pueden reconocerse a partir de fotografías, relatos y cartografía.

\subsection{Criterios de lectura}

El diagnóstico se organiza a partir de cinco criterios observables:

\begin{enumerate}
	\item \textbf{Sistema constructivo}: materiales principales, disposición de muros, tipo de cerramiento y ensambles visibles.
	\item \textbf{Uso actual o reciente}: función doméstica, productiva o de resguardo atribuida por los habitantes.
	\item \textbf{Estado físico}: permanencias, deterioros, sustituciones y elementos agregados.
	\item \textbf{Relación con el entorno}: proximidad de materiales, vegetación, caminos, patios o áreas productivas.
	\item \textbf{Relato asociado}: memoria de construcción, mantenimiento, modificación o abandono.
\end{enumerate}

La lectura se centra en la relación entre materialidad, territorio y saber local, no únicamente en la forma visible de cada construcción.

\subsection{Mapas y territorio}

La cartografía se plantea en una secuencia de escalas. Los planos PL-01, PL-02 y PL-03 ubican el caso desde el contexto nacional y estatal hasta el municipio de El Cardonal y las comunidades de El Deca y El Buena. El plano PL-04 registra las construcciones tradicionales visualizadas en campo y distingue los cuatro casos seleccionados para el diagnóstico. El plano PL-05 relaciona esos cuatro sistemas con relieve, vegetación, caminos y perímetros comunitarios, sin convertir el mapa en una prueba de extracción exacta de materiales. Finalmente, los planos PL-06 y PL-07 muestran los recorridos aproximados de documentación realizados en mayo de 2025 y el área general donde se concentró el registro constructivo.

Los mapas no se presentan como prueba aislada, sino como una forma de relacionar las construcciones con su contexto: caminos, zonas de habitación, áreas de cultivo, presencia de vegetación útil y posibles recorridos de acceso a materiales. Esta lectura permite reforzar la hipótesis de que los sistemas constructivos responden a una disponibilidad material localizada y a decisiones domésticas acumuladas. Los recorridos se representan como trazos aproximados sobre caminos de OpenStreetMap, porque esta base muestra accesos locales que no siempre aparecen con el mismo detalle en la cartografía institucional. En la versión pública de la ICR, la ubicación se maneja por comunidad, por clave de registro y por área aproximada, no mediante coordenadas exactas.

La lectura cartográfica también distingue entre el perímetro oficial de localidad y el territorio reconocido por sus habitantes. En el registro, los casos C-02 y C-04, junto con una construcción adicional no seleccionada, se ubican fuera del polígono geoestadístico de El Deca; sin embargo, las personas propietarias los reconocen como pertenecientes a esa comunidad. Por ello, el diagnóstico no corrige la adscripción comunitaria a partir del límite cartográfico, sino que registra la tensión entre una delimitación institucional y una pertenencia comunitaria sostenida por relaciones familiares, uso cotidiano, recorridos y memoria local \citep{ziccardi2015habitabilidad, cortes2014juventud}.

Esta diferencia también aparece en la red de caminos representada. La capa de vialidades empleada para el mapa local corresponde a ejes de vialidad del Marco Geoestadístico, por lo que no necesariamente incluye todos los accesos rurales, veredas o terracerías utilizados localmente. Al revisar las capas lineales complementarias de INEGI, el acceso más cercano a C-02 y C-04 aparece como camino de terracería hacia El Deca, mientras que la base de OpenStreetMap muestra trazos de acceso más próximos. Por tanto, los mapas deben leerse como una triangulación entre cartografía oficial, cartografía colaborativa y relato local, no como una representación exhaustiva de los recorridos comunitarios.

\begin{figure}[p]
	\centering
	\mapaicr{PL-01_Contexto_Mexico_Hidalgo.png}
	\caption{Mapa PL-01. Contexto de México e Hidalgo.}
	\label{fig:pl01_contexto_mexico_hidalgo}
\end{figure}

\begin{figure}[p]
	\centering
	\mapaicr{PL-02_Hidalgo_Cardonal.png}
	\caption{Mapa PL-02. Ubicación de Cardonal dentro de Hidalgo y el Valle del Mezquital.}
	\label{fig:pl02_hidalgo_cardonal}
\end{figure}

\begin{figure}[p]
	\centering
	\mapaicr{PL-03_Cardonal_Comunidades.png}
	\caption{Mapa PL-03. Comunidades de El Deca y El Buena dentro de Cardonal.}
	\label{fig:pl03_cardonal_comunidades}
\end{figure}

\begin{figure}[p]
	\centering
	\mapaicr{PL-04_Construcciones_Registradas.png}
	\caption{Mapa PL-04. Construcciones tradicionales registradas y cuatro casos seleccionados para la ICR.}
	\label{fig:pl04_construcciones_registradas}
\end{figure}

\begin{figure}[p]
	\centering
	\mapaicr{PL-05_Relacion_Materiales_Territorio.png}
	\caption{Mapa PL-05. Relación entre sistemas constructivos, relieve, vegetación, caminos y comunidades.}
	\label{fig:pl05_materiales_territorio}
\end{figure}

\begin{figure}[p]
	\centering
	\mapaicr{PL-06_Recorridos_Documentacion_2025.png}
	\caption{Mapa PL-06. Recorridos aproximados de documentación realizados en mayo de 2025.}
	\label{fig:pl06_recorridos_documentacion}
\end{figure}

\begin{figure}[p]
	\centering
	\mapaicr{PL-07_Area_Registro_Constructivo.png}
	\caption{Mapa PL-07. Área aproximada donde se concentró el registro de construcciones tradicionales.}
	\label{fig:pl07_area_registro_constructivo}
\end{figure}

\FloatBarrier

\subsection{Fichas por caso}

Cada construcción se presenta mediante una ficha sintética. La ficha incluye:

\begin{itemize}
	\item clave de identificación anónima;
	\item comunidad;
	\item sistema constructivo predominante;
	\item materiales observados;
	\item uso actual o uso recordado;
	\item estado físico cualitativo;
	\item cambios visibles o relatados;
	\item evidencia fotográfica;
	\item observaciones cartográficas;
	\item fragmento o síntesis del relato asociado.
\end{itemize}

No se emplea un código único de estado físico, porque ese recurso simplificaría en exceso la condición material, morfológica, social y narrativa de cada construcción. En su lugar, se describe cada caso atendiendo permanencias, adecuaciones, sustituciones, uso y relato asociado.

\subsubsection{Ficha C-01: construcción con penca de maguey}

El caso C-01 se ubica públicamente en la comunidad de El Buena y corresponde a una construcción de penca de maguey usada como bodega vinculada a la recolección, venta y almacenamiento de pulque. Su materialidad principal está compuesta por pencas de maguey, horcones, quiote de maguey, quiote de lechuguilla, pilotes y tablones de madera. La construcción se mantiene en uso y conserva el sistema constructivo principal, aunque presenta adecuaciones materiales puntuales.

Los cambios más visibles son la sustitución de la cubierta originalmente asociada a paja o penca de maguey por lámina de cartón y la incorporación de un piso de concreto. Estos cambios no eliminan la forma general ni el cerramiento principal de penca de maguey, pero sí muestran una adaptación funcional del espacio. Según el relato, el piso fue incorporado porque los tambos usados para guardar el pulque absorbían humedad y se deterioraban con rapidez.

El relato asociado vincula la construcción con la continuidad familiar del raspado de pulque. El informante señaló que antes existían más magueyes en sus magueyales, pero que el consumo local de pulque disminuyó frente a la preferencia por la cerveza. También explicó que las pencas deben reponerse cuando se pudren y que, de acuerdo con enseñanzas familiares, conviene cortarlas después de la luna llena para que duren más tiempo.

\begin{figure}[!htbp]
	\centering
	\setlength{\tabcolsep}{2pt}
	\begin{tabular}{cc}
		\fotocaso{C-01_penca_maguey_publicable_01.jpg} & \fotocaso{C-01_penca_maguey_publicable_02.jpg} \\
		\fotocaso{C-01_penca_maguey_publicable_03.jpg} & \fotocaso{C-01_penca_maguey_publicable_04.jpg} \\
		\fotocaso{C-01_penca_maguey_publicable_05.jpg} & \fotocaso{C-01_penca_maguey_publicable_06.jpg}
	\end{tabular}
	\caption{Registro fotográfico publicable del caso C-01, construcción con penca de maguey, primera lámina.}
	\label{fig:c01_fotos_publicables_01}
\end{figure}

\begin{figure}[!htbp]
	\centering
	\setlength{\tabcolsep}{2pt}
	\begin{tabular}{cc}
		\fotocasogrande{C-01_penca_maguey_publicable_07.jpg} & \fotocasogrande{C-01_penca_maguey_publicable_08.jpg} \\
		\fotocasogrande{C-01_penca_maguey_publicable_09.jpg} & \fotocasogrande{C-01_penca_maguey_publicable_10.jpg} \\
		\fotocasogrande{C-01_penca_maguey_publicable_11.jpg} & \fotocasogrande{C-01_penca_maguey_publicable_12.jpg}
	\end{tabular}
	\caption{Registro fotográfico publicable del caso C-01, construcción con penca de maguey, segunda lámina.}
	\label{fig:c01_fotos_publicables_02}
\end{figure}

\begin{figure}[!htbp]
	\centering
	\setlength{\tabcolsep}{2pt}
	\begin{tabular}{cc}
		\fotocasomuygrande{C-01_penca_maguey_publicable_13.jpg} & \fotocasomuygrande{C-01_penca_maguey_publicable_14.jpg} \\
		\multicolumn{2}{c}{\fotocasocentrada{C-01_penca_maguey_publicable_15.jpg}}
	\end{tabular}
	\caption{Registro fotográfico publicable del caso C-01, construcción con penca de maguey, tercera lámina.}
	\label{fig:c01_fotos_publicables_03}
\end{figure}

\FloatBarrier

\subsubsection{Ficha C-02: construcción de piedra local o ``tepetate''}

El caso C-02 se registra como perteneciente a El Deca por adscripción local, aunque queda fuera de los perímetros geoestadísticos oficiales de El Deca y El Buena. La construcción corresponde a un sistema de piedra local nombrada por la comunidad como ``tepetate''. Esta denominación se conserva porque forma parte del conocimiento local; sin embargo, no se afirma una clasificación petrográfica, ya que no se realizó un análisis especializado del material.

La construcción fue habitada por la familia durante varios años y posteriormente quedó como bodega de herramienta y alimento para ganado, con un entorno usado como corral para cabras y magueyes. Se mantiene en pie, con permanencia morfológica general, pero presenta adecuaciones materiales importantes: sustitución de la cubierta por lámina de acero y recubrimiento con cemento en gran parte del muro. Todavía se conservan sectores donde la piedra se encuentra asentada con lodo.

El relato permite entender que esas adecuaciones no fueron hechas solo por preferencia formal, sino como respuesta a un problema cotidiano. La familia señaló que los roedores podían entrar en los muros porque la mezcla de lodo era suave y permitía la formación de huecos. Por ello se cubrieron partes del muro con cemento.

\begin{figure}[!h]
	\centering
	\setlength{\tabcolsep}{2pt}
	\begin{tabular}{cc}
		\fotocaso{C-02_piedra_local_tepetate_publicable_01.jpg} & \fotocaso{C-02_piedra_local_tepetate_publicable_02.jpg}
	\end{tabular}
	\caption{Registro fotográfico publicable del caso C-02, construcción de piedra local / tepetate, primera lámina.}
	\label{fig:c02_fotos_publicables_01}
\end{figure}

\begin{figure}[!htbp]
	\centering
	\setlength{\tabcolsep}{2pt}
	\begin{tabular}{cc}
		\fotocasogrande{C-02_piedra_local_tepetate_publicable_03.jpg} & \fotocasogrande{C-02_piedra_local_tepetate_publicable_04.jpg} \\
		\fotocasogrande{C-02_piedra_local_tepetate_publicable_05.jpg} & \fotocasogrande{C-02_piedra_local_tepetate_publicable_06.jpg}
	\end{tabular}
	\caption{Registro fotográfico publicable del caso C-02, construcción de piedra local / tepetate, segunda lámina.}
	\label{fig:c02_fotos_publicables_02}
\end{figure}

\begin{figure}[!htbp]
	\centering
	\setlength{\tabcolsep}{2pt}
	\begin{tabular}{cc}
		\fotocasogrande{C-02_piedra_local_tepetate_publicable_07.jpg} & \fotocasogrande{C-02_piedra_local_tepetate_publicable_08.jpg} \\
		\fotocasogrande{C-02_piedra_local_tepetate_publicable_09.jpg} & \fotocasogrande{C-02_piedra_local_tepetate_publicable_10.jpg} \\
		\fotocasogrande{C-02_piedra_local_tepetate_publicable_11.jpg} & \fotocasogrande{C-02_piedra_local_tepetate_publicable_12.jpg}
	\end{tabular}
	\caption{Registro fotográfico publicable del caso C-02, construcción de piedra local / tepetate, tercera lámina.}
	\label{fig:c02_fotos_publicables_03}
\end{figure}

\begin{figure}[!htbp]
	\centering
	\setlength{\tabcolsep}{2pt}
	\begin{tabular}{cc}
		\fotocasogrande{C-02_piedra_local_tepetate_publicable_13.jpg} & \fotocasogrande{C-02_piedra_local_tepetate_publicable_14.jpg} \\
		\fotocasogrande{C-02_piedra_local_tepetate_publicable_15.jpg} & \fotocasogrande{C-02_piedra_local_tepetate_publicable_16.jpg} \\
		\multicolumn{2}{c}{\fotocasocentrada{C-02_piedra_local_tepetate_publicable_17.jpg}}
	\end{tabular}
	\caption{Registro fotográfico publicable del caso C-02, construcción de piedra local / tepetate, cuarta lámina.}
	\label{fig:c02_fotos_publicables_04}
\end{figure}

\FloatBarrier

\subsubsection{Ficha C-03: construcción de bajareque}

El caso C-03 se ubica en El Buena y corresponde a una cocina de bajareque a base de fogón. Es el caso con mayor permanencia morfológica dentro del corpus: conserva la forma general, el entramado y el sistema constructivo principal. Sus materiales incluyen quiote de lechuguilla, quiote de maguey, castillo (\emph{Leonotis nepetifolia}), polines de madera, tierra, lámina de cartón y horcones.

El principal cambio material registrado es la sustitución de la cubierta por lámina de cartón. A diferencia de otros casos, no se identificó una alteración dominante del sistema de muros. La importancia de esta construcción se relaciona tanto con su materialidad como con su uso cotidiano: continúa funcionando como cocina de fogón.

El relato de una mujer mayor de la casa muestra un fuerte apego afectivo al espacio. Ella señaló que en las estufas nuevas la comida no sabe igual y que conserva la cocina porque la asocia con el sabor de su comida y con una forma propia de cocinar. También expresó que, cuando ella fallezca, probablemente la cocina será demolida o sustituida. Esta posibilidad convierte al caso en un ejemplo de riesgo generacional: la amenaza no proviene solo del deterioro material, sino de la pérdida de la persona que sostiene el uso, el afecto y el sentido del espacio. El relato también explica el uso del castillo como relleno, debido a su disponibilidad en el jardín.

\begin{figure}[!h]
	\centering
	\setlength{\tabcolsep}{2pt}
	\begin{tabular}{cc}
		\fotocaso{C-03_bajareque_publicable_01.jpg} & \fotocaso{C-03_bajareque_publicable_03.jpg}
	\end{tabular}
	\caption{Registro fotográfico publicable del caso C-03, construcción de bajareque, primera lámina.}
	\label{fig:c03_fotos_publicables_01}
\end{figure}

\begin{figure}[!htbp]
	\centering
	\setlength{\tabcolsep}{2pt}
	\begin{tabular}{cc}
		\fotocaso{C-03_bajareque_publicable_05.jpg} & \fotocaso{C-03_bajareque_publicable_07.jpg} \\
		\fotocaso{C-03_bajareque_publicable_08.jpg} & \fotocaso{C-03_bajareque_publicable_09.jpg} \\
		\fotocaso{C-03_bajareque_publicable_10.jpg} & \fotocaso{C-03_bajareque_publicable_11.jpg}
	\end{tabular}
	\caption{Registro fotográfico publicable del caso C-03, construcción de bajareque, segunda lámina.}
	\label{fig:c03_fotos_publicables_02}
\end{figure}

\begin{figure}[!htbp]
	\centering
	\setlength{\tabcolsep}{2pt}
	\begin{tabular}{cc}
		\fotocaso{C-03_bajareque_publicable_12.jpg} & \fotocaso{C-03_bajareque_publicable_13.jpg} \\
		\fotocaso{C-03_bajareque_publicable_14.jpg} & \fotocaso{C-03_bajareque_publicable_15.jpg} \\
		\fotocaso{C-03_bajareque_publicable_16.jpg} & \fotocaso{C-03_bajareque_publicable_17.jpg}
	\end{tabular}
	\caption{Registro fotográfico publicable del caso C-03, construcción de bajareque, tercera lámina.}
	\label{fig:c03_fotos_publicables_03}
\end{figure}

\begin{figure}[!htbp]
	\centering
	\setlength{\tabcolsep}{2pt}
	\begin{tabular}{cc}
		\fotocaso{C-03_bajareque_publicable_18.jpg} & \fotocaso{C-03_bajareque_publicable_19.jpg} \\
		\fotocaso{C-03_bajareque_publicable_20.jpg} & \fotocaso{C-03_bajareque_publicable_21.jpg} \\
		\fotocaso{C-03_bajareque_publicable_22.jpg} & \fotocaso{C-03_bajareque_publicable_23.jpg} \\
		\multicolumn{2}{c}{\fotocasotexto{C-03_bajareque_publicable_24.jpg}}
	\end{tabular}
	\caption{Registro fotográfico publicable del caso C-03, construcción de bajareque, cuarta lámina.}
	\label{fig:c03_fotos_publicables_04}
\end{figure}

\FloatBarrier

\subsubsection{Ficha C-04: cerramiento de órganos}

El caso C-04 se registra como perteneciente a El Deca por adscripción local y, al igual que C-02, queda fuera de los perímetros geoestadísticos oficiales de El Deca y El Buena. Corresponde a una cocina a base de fogón y almacenamiento de frutas y verduras, delimitada por cactus órgano (\emph{Pachycereus marginatus}). Este caso es relevante porque el órgano suele emplearse como barrera divisoria en distintas partes del paisaje doméstico, pero aquí funciona como cerramiento de un espacio específico de uso cotidiano.

La construcción se encuentra en uso y conserva la mayor parte del cerramiento vegetal. Sin embargo, uno de sus muros fue parcialmente sustituido por láminas de diferentes materiales debido a la muerte de algunos órganos. La cubierta es de lámina de cartón. Por ello, el estado físico no puede describirse únicamente como conservado o deteriorado: combina permanencia del cerramiento vivo o semivivo, sustitución localizada y continuidad de uso.

El relato de la dueña destaca que la cocina, aunque pequeña, resulta cómoda y adecuada para sus necesidades. También explica un conocimiento específico sobre el manejo del material: los órganos deben trasplantarse después de la luna llena para adaptarse mejor al nuevo sitio, y se debe elegir un tamaño adecuado para que mantengan proporciones semejantes entre sí. La abundancia de órganos en la zona facilita la selección de ejemplares.

\begin{figure}[!h]
	\centering
	\setlength{\tabcolsep}{2pt}
	\begin{tabular}{cc}
		\fotocaso{C-04_organos_publicable_01.jpg} & \fotocaso{C-04_organos_publicable_03.jpg}
	\end{tabular}
	\caption{Registro fotográfico publicable del caso C-04, cerramiento de órganos, primera lámina.}
	\label{fig:c04_fotos_publicables_01}
\end{figure}

\begin{figure}[!htbp]
	\centering
	\setlength{\tabcolsep}{2pt}
	\begin{tabular}{cc}
		\fotocaso{C-04_organos_publicable_05.jpg} & \fotocaso{C-04_organos_publicable_07.jpg} \\
		\fotocaso{C-04_organos_publicable_08.jpg} & \fotocaso{C-04_organos_publicable_09.jpg} \\
		\fotocaso{C-04_organos_publicable_10.jpg} & \fotocaso{C-04_organos_publicable_12.jpg}
	\end{tabular}
	\caption{Registro fotográfico publicable del caso C-04, cerramiento de órganos, segunda lámina.}
	\label{fig:c04_fotos_publicables_02}
\end{figure}

\begin{figure}[!htbp]
	\centering
	\setlength{\tabcolsep}{2pt}
	\begin{tabular}{cc}
		\fotocaso{C-04_organos_publicable_13.jpg} & \fotocaso{C-04_organos_publicable_14.jpg} \\
		\fotocaso{C-04_organos_publicable_15.jpg} & \fotocaso{C-04_organos_publicable_16.jpg} \\
		\multicolumn{2}{c}{\fotocasotexto{C-04_organos_publicable_17.jpg}}
	\end{tabular}
	\caption{Registro fotográfico publicable del caso C-04, cerramiento de órganos, tercera lámina.}
	\label{fig:c04_fotos_publicables_03}
\end{figure}

\FloatBarrier

\subsection{Comparación entre sistemas}

Después de presentar las fichas individuales, el análisis compara los sistemas constructivos según sus componentes principales. La comparación permite identificar patrones como:

\begin{itemize}
	\item uso de materiales vegetales frente a materiales pétreos o térreos;
	\item presencia de cubiertas sustituidas por lámina;
	\item permanencia morfológica aun cuando existen sustituciones materiales;
	\item relación entre el sistema constructivo y el uso actual de la construcción;
	\item cambios asociados a mantenimiento, adaptación o sustitución parcial.
\end{itemize}

La comparación no busca establecer una jerarquía de valor entre sistemas, sino describir su diversidad y las condiciones que explican su permanencia o transformación.

Los cuatro casos muestran que las construcciones tradicionales no permanecen por ausencia de cambio, sino por una combinación de continuidad morfológica, adecuaciones materiales y uso cotidiano. En todos aparecen sustituciones en cubierta, principalmente con lámina, pero estas no eliminan por sí solas la lectura del sistema cuando se conserva la forma general, el cerramiento principal o la lógica material del espacio.

La relación con el territorio se expresa de maneras distintas. En C-01, el maguey articula construcción, mantenimiento y actividad pulquera. En C-02, la piedra local denominada tepetate por la comunidad vincula memoria habitacional, uso actual y adaptación frente a roedores. En C-03, el bajareque combina tierra, quiotes y castillo disponible en el jardín con un valor doméstico asociado al sabor y al fogón. En C-04, los órganos funcionan como cerramiento vivo o semivivo de una cocina, no solo como barrera divisoria.

El estado físico no se reduce a códigos simples. Cada caso requiere una descripción cualitativa que distinga forma, material, uso, mantenimiento, relato y riesgo de continuidad. Las implicaciones de conservación se orientan hacia documentación, prevención y transmisión de conocimiento, más que hacia una restauración arquitectónica convencional.

\subsection{Síntesis operativa del diagnóstico}

El diagnóstico cierra con una síntesis operativa que servirá como puente entre el análisis académico, las fichas y los criterios de documentación:

\begin{itemize}
	\item \textbf{C-01} representa mantenimiento activo de un sistema asociado al pulque y al maguey. Su conservación depende de la reposición de pencas, del conocimiento sobre el corte y de la continuidad de una práctica productiva.
	\item \textbf{C-02} representa permanencia morfológica con adecuaciones materiales derivadas de problemas cotidianos de uso. Debe nombrarse como piedra local / tepetate, respetando el término comunitario sin afirmar una clasificación petrográfica.
	\item \textbf{C-03} representa la mayor permanencia formal y un alto valor afectivo, pero también un riesgo generacional asociado a la posible pérdida de la usuaria principal.
	\item \textbf{C-04} representa un cerramiento vegetal vivo o semivivo aplicado a una cocina, con conocimiento específico de selección, trasplante, tamaño y reposición de órganos.
\end{itemize}
